\section{Background and Related Work}
\label{sec:background}

This section reviews the mathematical foundations of the algorithms under
study.  We begin with the classical Particle Swarm Optimisation framework and
its principal refinements, proceed to the Quantum-Behaved extension and the
Weighted Mean Best improvement, and conclude with a brief survey of the
benchmark functions employed in our experimental evaluation.

%% ----------------------------------------------------------------
\subsection{Classical Particle Swarm Optimisation}
\label{sec:bg:pso}

Particle Swarm Optimisation, introduced by Kennedy and Eberhart in
1995~\cite{Kennedy1995}, models the collective behaviour of a swarm of
$N$ particles moving through a $D$-dimensional search space.  Each particle
$i$ maintains a position vector $\mathbf{x}_i \in \mathbb{R}^D$ and a
velocity vector $\mathbf{v}_i \in \mathbb{R}^D$.  At every iteration $t$, the
velocity is updated according to the cognitive--social rule and the position is
advanced accordingly:
\begin{align}
  \mathbf{v}_i^{t+1} &= w\,\mathbf{v}_i^{t}
    + c_1\,\mathbf{r}_1 \odot \bigl(\mathbf{pbest}_i - \mathbf{x}_i^{t}\bigr)
    + c_2\,\mathbf{r}_2 \odot \bigl(\mathbf{gbest} - \mathbf{x}_i^{t}\bigr),
  \label{eq:pso_velocity} \\[4pt]
  \mathbf{x}_i^{t+1} &= \mathbf{x}_i^{t} + \mathbf{v}_i^{t+1},
  \label{eq:pso_position}
\end{align}
where $\mathbf{pbest}_i$ denotes the best position visited by particle~$i$
since initialisation, $\mathbf{gbest}$ denotes the best position found by any
particle in the swarm, $c_1$ and $c_2$ are the cognitive and social
acceleration coefficients, $\mathbf{r}_1$ and $\mathbf{r}_2$ are vectors of
independent uniform random variables drawn from $\mathcal{U}(0,1)$, and
$\odot$ represents the Hadamard (element-wise) product.  The scalar $w$ is the
\emph{inertia weight}, introduced by Shi and Eberhart~\cite{Shi1998} to
control the influence of the particle's previous velocity on its current
trajectory.

The inertia weight plays a critical role in balancing exploration and
exploitation.  A large value of $w$ encourages broad search by preserving the
particle's momentum, while a small value promotes local refinement by damping
velocity and allowing the cognitive and social terms to dominate.  In practice,
$w$ is commonly decayed linearly from an initial value $w_{\max}$ to a final
value $w_{\min}$ over the course of the optimisation:
\begin{equation}
  w(t) = w_{\max} - \frac{t}{T}\bigl(w_{\max} - w_{\min}\bigr),
  \label{eq:inertia_decay}
\end{equation}
where $T$ is the total number of iterations.  Typical choices are
$w_{\max} = 0.9$ and $w_{\min} = 0.4$, values that have been found to provide
a good balance across a variety of benchmark landscapes.

Clerc and Kennedy provided an alternative perspective through the
\emph{constriction coefficient} approach~\cite{Clerc2002}.  By analysing the
eigenvalues of the recurrence relation governing particle trajectories, they
showed that convergence to a fixed point is guaranteed when the sum of the
acceleration coefficients exceeds a critical threshold, $c_1 + c_2 > 4$, and a
constriction factor $\chi$ is applied:
\begin{equation}
  \chi = \frac{2}{\bigl|2 - \varphi - \sqrt{\varphi^2 - 4\varphi}\,\bigr|},
  \qquad \varphi = c_1 + c_2 > 4.
  \label{eq:constriction}
\end{equation}
The velocity update then becomes
$\mathbf{v}_i^{t+1} = \chi\bigl[\mathbf{v}_i^{t} + c_1\,\mathbf{r}_1 \odot
(\mathbf{pbest}_i - \mathbf{x}_i^{t}) + c_2\,\mathbf{r}_2 \odot
(\mathbf{gbest} - \mathbf{x}_i^{t})\bigr]$.  With $c_1 = c_2 = 2.05$ one
obtains $\chi \approx 0.7298$, which yields a net cognitive and social pull
equivalent to $c_1' = c_2' \approx 1.494$---precisely the parameterisation
used in the present study.

To prevent particles from acquiring unbounded velocities during the early
exploration phase, a velocity clamping mechanism is applied.  The velocity in
each dimension is clipped to the interval $[-v_{\max},\; v_{\max}]$, where
$v_{\max}$ is typically set to a fraction of the search range.  This
safeguard does not alter the convergence properties established by the
constriction analysis but improves numerical stability in practice.


%% ----------------------------------------------------------------
\subsection{Quantum-Behaved Particle Swarm Optimisation}
\label{sec:bg:qpso}

Classical PSO updates particle positions through deterministic velocity
increments, a mechanism that confines each particle to a trajectory governed by
its current velocity and the attractor landscape defined by personal and global
bests.  Sun, Feng, and Xu observed that this trajectory-based formulation
prevents particles from escaping deep local basins unless stochastic
perturbations are large enough to vault over the intervening
barrier~\cite{Sun2004}.  Drawing on the quantum-mechanical concept of a
particle bound in a potential well---where the wave function extends beyond the
classical turning points with exponentially decaying amplitude---they proposed
QPSO, a variant in which particles have no velocity and their positions are
instead sampled from a probability distribution that inherently permits
\emph{tunnelling} through potential barriers.

The derivation proceeds from a simplified quantum model.  Consider a particle
moving in a one-dimensional delta potential well centred at an attractor
point~$p$.  The time-independent Schr\"odinger equation yields a ground-state
wave function
$\psi(x) \propto \exp\bigl(-|x - p|/L\bigr)$,
whose squared modulus gives the position probability density
$|\psi(x)|^2 \propto \exp\bigl(-2\,|x - p|/L\bigr)$.
The characteristic length $L$ controls the spatial extent of the distribution:
a large $L$ produces a broad, exploratory distribution, while a small $L$
yields a narrow, exploitative one.  Sampling from this double-exponential
distribution via inverse transform sampling gives the position update rule
that forms the core of QPSO.

In the multi-dimensional formulation, the attractor point for particle~$i$ at
iteration~$t$ is computed as a random interpolation between the particle's
personal best and the swarm's global best:
\begin{equation}
  \mathbf{p}_i = \boldsymbol{\phi} \odot \mathbf{pbest}_i
    + (\mathbf{1} - \boldsymbol{\phi}) \odot \mathbf{gbest},
  \label{eq:qpso_attractor}
\end{equation}
where $\boldsymbol{\phi}$ is a vector of independent uniform random variables
on $[0,1]$.  The \emph{mean best position} (mbest) is defined as the centroid
of all personal bests in the swarm:
\begin{equation}
  \mathbf{mbest} = \frac{1}{N} \sum_{j=1}^{N} \mathbf{pbest}_j.
  \label{eq:mbest}
\end{equation}
The quantum displacement for particle~$i$ is then
\begin{equation}
  \boldsymbol{\delta}_i = \beta\,
    \bigl|\mathbf{mbest} - \mathbf{x}_i^{t}\bigr|
    \odot \ln\!\Bigl(\frac{1}{\mathbf{u}}\Bigr),
  \label{eq:qpso_displacement}
\end{equation}
where $\beta$ is the \emph{contraction--expansion coefficient} and
$\mathbf{u}$ is a vector of independent uniform random variables on $(0,1)$.
The logarithmic term arises from the inverse CDF of the double-exponential
distribution, and the factor $|\mathbf{mbest} - \mathbf{x}_i^{t}|$ scales the
step size in proportion to the particle's distance from the swarm's consensus
position, ensuring that outlying particles explore more broadly while
particles near the consensus region exploit more intensively.

The position is updated by adding or subtracting the displacement from the
attractor:
\begin{equation}
  \mathbf{x}_i^{t+1} = \mathbf{p}_i \pm \boldsymbol{\delta}_i,
  \label{eq:qpso_update}
\end{equation}
where the sign is chosen independently for each dimension with equal
probability.  This binary sign decision constitutes the \emph{tunnelling
direction}: a positive sign moves the particle outward from the attractor,
while a negative sign moves it inward.  In our framework, this sign is the
quantity most naturally suited to generation by a quantum circuit, as it is an
inherently discrete, binary random variable.

The contraction--expansion coefficient $\beta$ controls the width of the
quantum potential well and is analogous to the inertia weight in classical PSO.
Following the guidelines of Sun et al.~\cite{Sun2012}, $\beta$ is decayed
linearly from an initial value $\beta_{\max}$ (broad exploration) to a final
value $\beta_{\min}$ (tight exploitation) over the course of the optimisation:
\begin{equation}
  \beta(t) = \beta_{\max}
    - \frac{t}{T}\bigl(\beta_{\max} - \beta_{\min}\bigr).
  \label{eq:beta_decay}
\end{equation}
Sun et al.\ recommended $\beta_{\max} = 1.0$ and $\beta_{\min} = 0.5$ as
default values, and their analysis of individual particle behaviour confirmed
that this range ensures convergence while maintaining sufficient
diversity~\cite{Sun2012}.  In our implementation, we additionally widen $\beta$
by a factor of 1.3 when stagnation is detected, temporarily broadening the
quantum well to facilitate escape from local basins before resuming the
standard decay schedule.


%% ----------------------------------------------------------------
\subsection{Weighted Mean Best Position (WQPSO)}
\label{sec:bg:wqpso}

The standard mean best position defined in Equation~\eqref{eq:mbest} assigns
equal weight to every particle's personal best, regardless of its fitness.  On
valley-type landscapes such as the Rosenbrock function, this can be
problematic: particles that have settled on the valley floor and those that
remain on the surrounding plateau contribute equally to $\mathbf{mbest}$,
pulling the consensus reference away from the narrow region of high-quality
solutions.

Xi, Sun, and Xu proposed an improved variant, Weighted QPSO
(WQPSO)~\cite{Xi2008}, in which each particle's contribution to the mean best
position is weighted by the inverse of its fitness.  Specifically, the
fitness-weighted mean best is defined as
\begin{equation}
  \mathbf{mbest}_w = \sum_{j=1}^{N} w_j\,\mathbf{pbest}_j,
  \qquad
  w_j = \frac{f_{\max} - f_j + \epsilon\,(f_{\max} - f_{\min})}
             {\displaystyle\sum_{k=1}^{N}
              \bigl[f_{\max} - f_k + \epsilon\,(f_{\max} - f_{\min})\bigr]},
  \label{eq:wqpso_mbest}
\end{equation}
where $f_j$ denotes the fitness of particle~$j$'s personal best, $f_{\max}$
and $f_{\min}$ are the worst and best fitness values in the swarm, and
$\epsilon$ is a small positive constant (set to 0.01 in our implementation)
that prevents any particle from receiving zero weight.  Under this scheme,
particles with superior fitness exert a stronger pull on the consensus
reference, biasing $\mathbf{mbest}_w$ toward the most promising region of the
search space.

The effect of fitness weighting is most pronounced on landscapes with narrow,
curved valleys.  On the Rosenbrock function, for instance, the global minimum
lies at the bottom of a parabolic valley that curves through the search space.
Particles that have descended into the valley possess substantially lower
fitness values than those that remain on the plateau, and the weighted mean
best consequently tracks the valley floor more closely than the unweighted
centroid.  This tighter alignment between the consensus reference and the
optimal region accelerates convergence along the valley without sacrificing
the broader exploratory behaviour governed by the contraction--expansion
coefficient~$\beta$.

In the implementation used in this work, QPSO always employs the
fitness-weighted mean best of Equation~\eqref{eq:wqpso_mbest}, while classical
PSO uses the standard unweighted formulation for its internal mbest tracking.
This choice reflects the finding of Xi et al.\ that the weighted variant
consistently matches or outperforms the original on the benchmark suite
considered, with no additional computational overhead beyond the weight
calculation itself.


%% ----------------------------------------------------------------
\subsection{Benchmark Functions}
\label{sec:bg:benchmarks}

To evaluate and compare optimisation algorithms, the community has developed a
standard collection of test functions whose landscape characteristics---number
of local optima, separability, conditioning, and valley structure---are well
understood.  The six functions employed in this study are summarised below and
in Table~\ref{tab:benchmarks}.

The \textbf{Sphere} function $f(\mathbf{x}) = \sum_{d=1}^{D} x_d^2$ is the
simplest benchmark: unimodal, separable, and symmetric about the origin.  Its
sole purpose is to verify that an algorithm converges at all and to establish a
baseline convergence rate.  Any competent swarm algorithm should solve the
Sphere function to machine precision within a modest iteration budget.

The \textbf{Rastrigin} function, introduced by Rastrigin in
1974~\cite{Rastrigin1974}, adds a cosine modulation to the Sphere:
$f(\mathbf{x}) = 10D + \sum_{d=1}^{D}\bigl[x_d^2 - 10\cos(2\pi x_d)\bigr]$.
This creates a regular grid of $\mathcal{O}(10^D)$ local minima separated by
barriers of height proportional to the cosine amplitude, making it one of the
most demanding tests of a swarm's ability to escape local basins.

The \textbf{Rosenbrock} function~\cite{Rosenbrock1960},
$f(\mathbf{x}) = \sum_{d=1}^{D-1}\bigl[100(x_{d+1} - x_d^2)^2
+ (1 - x_d)^2\bigr]$, is unimodal in two dimensions but features a narrow,
curved valley that is notoriously difficult for gradient-free methods to
navigate.  The global minimum at $\mathbf{x}^* = (1,\ldots,1)$ lies at the
bottom of this valley, and algorithms must track the valley floor over many
iterations to reach it.

The \textbf{Ackley} function~\cite{Ackley1987} combines an exponential term
involving the root-mean-square of the coordinates with a cosine term involving
the mean cosine:
$f(\mathbf{x}) = -20\exp\bigl(-0.2\sqrt{\frac{1}{D}\sum x_d^2}\bigr)
- \exp\bigl(\frac{1}{D}\sum\cos(2\pi x_d)\bigr) + 20 + e$.
Its landscape is characterised by a nearly flat outer region with many shallow
local minima surrounding a deep, narrow global basin at the origin, making it a
stringent test of an algorithm's exploitation capability.

The \textbf{Griewank} function~\cite{Griewank1981},
$f(\mathbf{x}) = \frac{1}{4000}\sum x_d^2
- \prod\cos\bigl(\frac{x_d}{\sqrt{d}}\bigr) + 1$,
produces a landscape with a large number of regularly distributed local
minima whose depth decreases with distance from the origin.  In low dimensions
the function is highly multimodal, but in higher dimensions the product term
dominates and the function becomes increasingly unimodal---an interesting
property that tests whether an algorithm's behaviour scales appropriately with
dimensionality.

The \textbf{Schwefel} function~\cite{Schwefel1981},
$f(\mathbf{x}) = 418.9829\,D - \sum x_d\sin\!\bigl(\sqrt{|x_d|}\bigr)$,
is particularly deceptive because the global minimum at
$\mathbf{x}^* \approx (420.97,\ldots,420.97)$ is geometrically distant from
the next-best local minima.  Algorithms that converge prematurely to the
nearest local basin have little chance of reaching the global optimum.

\begin{table}[t]
\centering
\caption{Benchmark functions used in the experimental evaluation.}
\label{tab:benchmarks}
\begin{tabular}{llccc}
\toprule
\textbf{Function} & \textbf{Landscape type} & \textbf{Bounds}
  & \textbf{$f(\mathbf{x}^*)$} & \textbf{$\mathbf{x}^*$} \\
\midrule
Sphere     & Unimodal, separable
  & $[-10,\; 10]^D$         & 0 & $\mathbf{0}$ \\
Rastrigin  & Multimodal, separable
  & $[-5.12,\; 5.12]^D$     & 0 & $\mathbf{0}$ \\
Rosenbrock & Valley, non-separable
  & $[-2.048,\; 2.048]^D$   & 0 & $\mathbf{1}$ \\
Ackley     & Multimodal, non-separable
  & $[-32.768,\; 32.768]^D$ & 0 & $\mathbf{0}$ \\
Griewank   & Multimodal, non-separable
  & $[-600,\; 600]^D$       & 0 & $\mathbf{0}$ \\
Schwefel   & Deceptive, separable
  & $[-500,\; 500]^D$       & 0 & $\approx 420.97$ \\
\bottomrule
\end{tabular}
\end{table}

Together, these six functions span a representative range of landscape
topologies.  The Sphere and Rosenbrock functions test convergence speed and
valley-following ability on unimodal landscapes.  The Rastrigin, Ackley, and
Griewank functions test escape from local optima under varying degrees of
multimodality.  The Schwefel function tests resilience against deceptive
gradients.  By evaluating all algorithm--mode combinations across this suite,
we obtain a comprehensive picture of relative performance that is unlikely to
be dominated by the idiosyncrasies of a single landscape.
